\begin{lemma}[Outside Kempe recolorings are stable]
Let $\varphi$ be a proper partial $k$-edge-coloring leaving exactly
$e_0$ uncolored.  If $\sigma$ is obtained by recoloring finitely many
full $(\alpha,\beta)$-chains contained in $G-V(T)$, then $\sigma$ is
proper, leaves exactly $e_0$ uncolored, and, for every control set
$C\subseteq[k]$, is $(T,C,\varphi)$-stable.  It is also
$(T,\varphi)$-invariant, and every exact edge identity incident with
$T$ has the same value under $\sigma$ and $\varphi$.
\end{lemma}
\begin{proof}
Let $\mathcal K$ be the selected family of full components from
\noderef{OutsideKempeRecoloring}.  We first prove that no selected edge
identity is incident with $T$.  If a selected edge-node $f$ were
incident with $v\in V(T)$, then the incidence edge in
\noderef{TwoColorIncidenceGraph} would join the vertex-node $v$ to
$f$.  The connected-component definition
\noderef{TwoColorComponent} would put $v$ in the same selected
component, contrary to the requirement that every vertex-node of each
selected component lies outside $V(T)$.  Therefore
\[
 \sigma(f)=\varphi(f)
 \quad\text{for every edge identity incident with }T.          \tag{1}
\]

Swapping two colors on a full two-color component preserves properness.
Indeed, at a vertex in a selected component, the incident
$\alpha$- and $\beta$-edges are exchanged together; at a vertex outside
that component no incident value changes.  Distinct
$(\alpha,\beta)$-components are vertex-disjoint by the connected
component relation in \noderef{TwoColorComponent}, so carrying out the
finitely many selected swaps produces no interaction between them.
All colors other than $\alpha,\beta$ are fixed.  Thus any equality of
new nonempty colors on adjacent distinct edge identities would pull back
to the same equality under $\varphi$, contradicting
\noderef{IsProperEdgeColoring}.  Hence $\sigma$ is proper.

The component swap fixes the value $\mathsf{none}$ and exchanges only
the two displayed nonempty values.  Consequently an edge is uncolored
under $\sigma$ exactly when it is uncolored under $\varphi$.
\noderef{IsOneEdgeUncolored} therefore shows that precisely $e_0$
remains uncolored.

Equation (1) holds for every edge incident with every $v\in V(T)$.
The incidence-witness definition of present colors, and then its
complement \noderef{MissingColors}, give
\[
 \overline\sigma(v)=\overline\varphi(v)\qquad(v\in V(T)).       \tag{2}
\]
For an arbitrary control set $C$, equation (1) is stronger than the
controlled edge clause of \noderef{StableColoring}, while (2) is its
vertex clause.  Thus $\sigma$ is $(T,C,\varphi)$-stable.
The displayed nonroot edges of $T$ are incident with $T$, so (1), along
with (2), is exactly \noderef{TreeColoringInvariant}.  This proves all
five conclusions.
\end{proof}
